\chapter{Case Study: "Interbolsa"}

\section{Executive Summary}

This case analyzes the collapse of InterBolsa S.A., the largest brokerage firm in Colombia in 2012, from the perspective of financial risk management and corporate governance. The study identifies how the combination of excessive leverage through repo operations, high concentration in an illiquid asset (Fabricato), structural dependence on short-term funding, and governance failures led to a systemic liquidity crisis.\\

The analysis demonstrates that internal risk models did generate meaningful early warning signals; however, those warnings were neither properly escalated nor incorporated into strategic decision-making. The progressive loss of confidence among banks, institutional investors, and market participants accelerated the deterioration until regulatory intervention became inevitable.

\vspace{0.5cm}

\section{Business Context}

During the first decade of the 21st century, Colombia experienced sustained growth in its capital markets, driven by economic expansion, increasing foreign investment, and institutional consolidation within the Colombian Stock Exchange (BVC).

Within this environment, InterBolsa became the leading brokerage firm in the country, expanding operations into Chile, Brazil, Panama, and the United States. Its business model evolved from traditional brokerage services toward sophisticated leveraged operations, repos, and financial structuring.

The firm operated under an organizational culture focused on aggressive growth and return maximization, where risk management gradually became subordinated to commercial objectives.

\vspace{0.5cm}

\section{Problem Statement}

The central problem of this case is to evaluate how a highly leveraged structure, financed through short-term repo operations and concentrated in a single illiquid collateral asset, generated a liquidity spiral and loss of confidence that ultimately resulted in InterBolsa's operational collapse.

This analysis seeks to answer the following questions:

\begin{itemize}
    \item Which risks were ignored or underestimated?
    \item How did concentration in Fabricato amplify systemic vulnerability?
    \item What governance and risk communication failures existed?
    \item Why were stress-testing results not incorporated into decision-making?
\end{itemize}

\vspace{0.5cm}

\section{Data Description}

\subsection{Sources of Information}

The analysis presented in this case is based on multiple sources of information, including the 2026 academic case study on InterBolsa, illustrative financial exhibits prepared for pedagogical purposes, and publicly available information from the Colombian Financial Superintendence (\textit{Superintendencia Financiera de Colombia}). Additionally, reconstructed datasets and simplified financial indicators were incorporated to replicate the structural dynamics of the crisis and facilitate the evaluation of liquidity, concentration, and systemic risk within the organization.

\subsection{Variables Analyzed}

\begin{table}[h!]
\centering
\begin{tabular}{|l|l|}
\hline
\textbf{Variable} & \textbf{Description} \\ \hline
Total Repos & Financial leverage level \\ \hline
Fabricato Concentration & Exposure to a single issuer \\ \hline
LCR Proxy & Liquidity coverage ratio \\ \hline
Haircuts & Market confidence deterioration \\ \hline
Bank Credit Lines & Dependence on external funding \\ \hline
Fabricato Price & Collateral deterioration \\ \hline
\end{tabular}
\caption{Key variables analyzed in the case}
\end{table}

\vspace{0.5cm}

\section{Methodology}

\subsection{Conceptual Framework}

This analysis applies principles of financial risk management related to \textbf{liquidity risk}, \textbf{concentration risk}, \textbf{market risk}, and \textbf{systemic risk}. Additionally, the study evaluates the role of \textbf{corporate governance} failures and the importance of \textbf{stress testing} frameworks in identifying vulnerabilities within highly leveraged financial institutions.

\vspace{0.3cm}

\subsection{Repo Financing Model}

Repo operations can be represented as:

\begin{equation}
\text{Obtained Liquidity} = \text{Collateral Value} \times (1 - h)
\end{equation}

where:

\begin{itemize}
    \item $h$ represents the haircut applied to collateral.
\end{itemize}

As Fabricato's price declined and haircuts increased, InterBolsa's refinancing capacity deteriorated rapidly.

\vspace{0.3cm}

\subsection{Stress Testing Framework}

The CRO developed stress scenarios based on potential \textbf{declines in collateral prices}, reductions in \textbf{bank funding lines}, increases in \textbf{margin calls}, and the simultaneous deterioration of \textbf{liquidity} and \textbf{market confidence}. These scenarios were designed to evaluate the firm's ability to withstand adverse market conditions and identify potential liquidity shortfalls under stressed financial environments.

The severe stress scenario assumed:

\begin{equation}
\Delta P_{Fabricato} \leq -50\%
\end{equation}

and bank funding availability reduced to:

\begin{equation}
L_{available} = 40\%
\end{equation}

The model projected a critical liquidity gap:

\begin{equation}
\text{Liquidity Gap} = -130 \text{ billion COP}
\end{equation}

\vspace{0.5cm}

\section{Case Analysis}

\subsection{Concentration Risk}

InterBolsa maintained approximately 56\% consolidated exposure to Fabricato, significantly exceeding internal prudential limits.

This concentration generated several vulnerabilities:

\begin{itemize}
    \item Excessive dependence on a single collateral asset
    \item Limited liquidation capacity in stressed markets
    \item High sensitivity to price fluctuations
    \item Reputational contagion risk
\end{itemize}

The progressive decline in Fabricato simultaneously deteriorated the \textbf{value of collateral}, weakened \textbf{creditor confidence}, and reduced the firm's \textbf{refinancing capacity}. As the market value of the collateral decreased, counterparties demanded additional guarantees and became increasingly reluctant to renew funding lines, accelerating the liquidity pressures faced by InterBolsa.

\vspace{0.3cm}

\subsection{Liquidity Risk}

The business model depended heavily on continuous refinancing through short-term repos.

Critical indicators included:

\begin{table}[h!]
\centering
\begin{tabular}{|l|c|}
\hline
\textbf{Indicator} & \textbf{October 2012} \\ \hline
LCR Proxy & 0.61x \\ \hline
Bank line utilization & 97\% \\ \hline
Average repo maturity & 18 days \\ \hline
Average haircut & 31\% \\ \hline
\end{tabular}
\caption{Critical liquidity indicators}
\end{table}

The loss of confidence triggered a classic liquidity crisis cascade. The initial decline in the \textbf{collateral price} generated significant \textbf{margin calls}, forcing InterBolsa to provide additional guarantees under increasingly adverse market conditions. As concerns regarding the firm's solvency intensified, banks began the \textbf{withdrawal of funding lines}, severely restricting access to short-term liquidity. To meet immediate obligations, the firm was forced into \textbf{distressed asset sales}, which further depressed market prices and amplified losses. This feedback loop ultimately evolved into a broader \textbf{systemic confidence crisis}, where counterparties, investors, and market participants simultaneously lost confidence in the institution's financial stability.

\vspace{0.3cm}

\subsection{Corporate Governance Failures}

One of the most important findings of the case was the information asymmetry between the risk management area and the Board of Directors.

While executive dashboards reported:

\begin{itemize}
    \item Diversified portfolio
    \item Satisfactory liquidity
    \item Controlled risk exposure
\end{itemize}

the technical risk department observed \textbf{critical concentration levels}, an \textbf{LCR below minimum regulatory and internal thresholds}, \textbf{excessive utilization of funding lines}, and \textbf{severe stress-test deficits}. These indicators collectively suggested that the firm's liquidity position was significantly more fragile than the executive reports presented to the Board of Directors, revealing substantial hidden vulnerabilities within the institution's funding structure.

This reflects a structural failure in corporate governance and organizational risk culture.

\vspace{0.3cm}

\subsection{Failures in Risk Management}

The main failures identified include:

\begin{enumerate}
    \item Underestimation of liquidity risk
    \item Ignoring stress-testing results
    \item Excessive dependence on short-term funding
    \item Incentives aligned with growth rather than risk control
    \item Ineffective escalation mechanisms for the CRO
    \item Overconfidence in the investment thesis
\end{enumerate}

\vspace{0.5cm}

\section{Results and Findings}

The case demonstrates that InterBolsa did not collapse solely because of financial losses, but rather as a consequence of a simultaneous crisis involving \textbf{liquidity deterioration}, the \textbf{loss of market confidence}, \textbf{weak corporate governance}, and \textbf{excessive concentration risk}. The interaction among these factors generated a self-reinforcing cycle in which declining liquidity conditions intensified market uncertainty, while governance failures prevented the organization from responding effectively to early warning signals and escalating systemic vulnerabilities.

Internal models identified the collapse scenario months in advance; however, the organization lacked both the governance structure and cultural willingness to act upon those warning signals.

\vspace{0.5cm}

\section{Risk Interpretation}

From an integrated risk management perspective, the InterBolsa case highlights:

\begin{itemize}
    \item The dangers of procyclical leverage
    \item The fragility of short-term funding models
    \item The critical importance of liquidity management
    \item The relevance of realistic stress testing
    \item The central role of confidence in financial markets
\end{itemize}

The case also illustrates that:

\begin{quote}
Liquidity disappears faster than solvency.
\end{quote}

Financial crises typically develop gradually over time and then materialize abruptly.

\vspace{0.5cm}

\section{Recommendations}

\begin{enumerate}
    \item Implement strict concentration limits
    \item Strengthen CRO independence
    \item Establish binding stress-testing frameworks
    \item Improve transparency toward the Board of Directors
    \item Diversify funding sources
    \item Develop formal liquidity contingency plans
    \item Incorporate forward-looking supervisory metrics
\end{enumerate}

\vspace{0.5cm}

\section{Limitations}

This analysis presents several limitations that should be considered when interpreting the results. First, the exhibits included in the case contain \textbf{illustrative and pedagogical data} rather than fully audited financial information. Additionally, complete \textbf{audited financial statements} and detailed internal records of InterBolsa were not publicly available, limiting the precision of certain quantitative assessments. Part of the analysis is therefore \textbf{reconstructive in nature}, relying on public reports, secondary sources, and simplified assumptions designed for academic discussion. Finally, the case was intentionally structured with a strong \textbf{pedagogical emphasis}, prioritizing conceptual understanding of liquidity crises, concentration risk, and corporate governance failures over forensic financial reconstruction.

\vspace{0.5cm}

\section{Conclusion}

The InterBolsa case represents one of the most important examples of systemic liquidity collapse in Latin American financial markets.

The firm did not fail exclusively because of losses associated with Fabricato, but because of the interaction between concentration risk, leverage, collateral deterioration, and the rapid disappearance of market confidence.

The principal lesson from this case is that quantitative models alone are insufficient if organizations lack a strong risk culture, effective governance mechanisms, and the willingness to act decisively on early warning signals.

\vspace{0.5cm}



\begin{thebibliography}{9}

\bibitem{Hull}
Hull, John.
\textit{Risk Management and Financial Institutions}.
Wiley.

\bibitem{McNeil}
McNeil, Frey \& Embrechts.
\textit{Quantitative Risk Management}.
Princeton University Press.

\bibitem{Basel}
Basel Committee on Banking Supervision.
\textit{Principles for Sound Liquidity Risk Management and Supervision}.

\bibitem{Interbolsa}
Academic Case Study:
\textit{INTERBOLSA: Rise, Expansion, and Collapse of Colombia's Largest Brokerage Firm}.
Financial and Enterprise Risk Management, 2026.

\end{thebibliography}
